\documentclass[a4paper, 12pt, oneside]{article}
\usepackage[utf8]{inputenc}
\usepackage[english]{babel}
\usepackage{fancyhdr}
\usepackage{geometry}
\usepackage{graphicx}
\usepackage{float}
\usepackage{subfig}
\usepackage{epstopdf}
\usepackage{amsmath}
\usepackage{amssymb}
\usepackage{siunitx}
\usepackage{enumitem}
\usepackage{booktabs}
\usepackage{gensymb}
\usepackage{makecell}
\usepackage{indentfirst}
\usepackage{bigints}
\usepackage{esint}
\usepackage{harpoon}
\usepackage{relsize}
\usepackage{CJKutf8}
\usepackage{cite}
\usepackage{url}
\usepackage{chemformula}
\usepackage{multicol}
\usepackage[useregional]{datetime2}
\usepackage[colorlinks = true,linkcolor = blue,hypertexnames = true]{hyperref}
\hypersetup{
    colorlinks=true,
    allcolors=blue	
}
\usepackage[super, square]{natbib}
\usepackage{multicol}
\usepackage{fancybox}
\usepackage{authblk}

\newcommand{\splitter}{\noindent\rule{\linewidth}{0.5pt}}

\geometry{a4paper, scale = 0.80}
\linespread{2.0}


\begin{document}
\begin{CJK}{UTF8}{bkai} % bsmi %bkai

\title{臺北市兜風指南}
\author[1]{張德權}
\affil[1]{Student at Department of Physics, National Taiwan University}






\date{\DTMdate{2021-12-12}}
\maketitle

\section*{關於兜風這件事}
坐上駕駛座，一條條四通八達的道路好像去除了人生的限制；目的地從不重要，享受的是這個不斷向前的過程，好似在追尋真理一般。不論是心情好或是心情差，我都喜歡兜風。馳騁在看不到盡頭的道路上，隨意的跟隨車流；暫時忘記心中的煩惱、拋開來自四面八方的壓力。既然兜風是如此的隨性，真的有出「兜風指南」的必要嗎？我想在台北市是有必要的。除非你會喜歡觀察路上每一個紅燈或是對他人的車尾特別有興趣，不然我相信「兜風指南」是會對你有幫助的。
\section*{時間選擇}
兜風是很看時間的，上下班尖峰時段肯定是不適合兜風的。擁擠的車潮佔據了馬路大部分的空間；右腳在油門與煞車間頻繁切換，與其說是兜風不如說是訓練耐心。因此，我推薦的兜風時間在清晨、下午以及 11：00pm 之後的時段；並且在上述時段中我最喜歡也最常選擇的就是 11：00pm 後的時段。
\section*{路線選擇}
臺北市是一個四通八達的城市，境內包高速公路、快速公路以及市區内快速道路。其中最不適合兜風的就屬高速公路了。撇開塞車不說，高速公路的車速本就讓人放鬆不下來，自然也就不適合兜風了。除去高速公路，臺北市最多的就是市區內快速道路了；這些路雖然俗稱快速道路，但實際上設計、罰則都適用於一般道路，因此超速罰單的金額也就微柔許多。
下面列出並介紹臺北市我認為適合兜風的道路：
\begin{itemize}
	\item 市民高架
	\item 環東大道
	\item 新生高架
	\item 環河南北快速
	\item 水源快速
	\item 國道三號甲線
	\item 敦化南北路
	\item 仁愛路
\end{itemize}
\subsection*{路線一：市民大道高架接環東大道}
市民大道是台北市中心最重要的東西向聯絡道，橫向貫穿臺北市中心，加上超高的高架高度，一覽城市風情就走這條。

環東大道西起八德路與基隆路的交會點，也就是是市民大道的終點，東至南港展覽館。多次經過總是覺得柏油品質很差，路面上的碎石多、很傷車頭烤漆，但是後段的雙層疊式高架給人城市進步的感覺，許多車評節目也會在此拍攝、取景。
\begin{itemize}
	\item 建議走法：大稻埕匝道上市民高架，接上環東大道。
	\item 建議速度：高速，體驗城市風情。
	\item 推薦指數：★★★☆☆
\end{itemize}
\subsection*{路線二：新生高架—環河南北快速—水源快速}
新生高架建於民國 73 年，屬於年代較為久遠的高架道路，但筆直的線型、偏少的匝道數量以及最重要的沒有一堆測速照相，使得這條高架道路總是給人舒服的感覺。

環河南北快速道路以及水源快速道路可以說是我最喜歡的兜風路線了。以師大水源路為界線往南是水源快，往北是環河快。這條路線沿著河岸，配合河堤的高度以及各個跨河橋樑，高度起起伏伏、高低錯落；讓人同時享受河岸風光以及好似開雲霄飛車的樂趣。
\begin{itemize}
	\item 建議走法：從濟南路匝道上新生高架，開到底後左轉劍潭路後左轉承德路；過了承德橋後右轉敦煌路，直走到底後左轉進入環河北路二段，直通水源快速道路；過了景美匝道後可以大腳油門衝一段，但請注意安全。
	\item 建議速度：中速，後段沿途多測速。
	\item 注意事項：水源快速道路有一個急彎，車速過快易出意外。
	\item 推薦指數：★★★★★
\end{itemize}
\subsection*{路線三：國道三號甲線}
編入國道，也是目前開唯一開放行駛大型重機的國道。連接辛亥路與木柵地區，沿途有兩個隧道，風景一般；整體沒有特別凸出，入選原因是離臺大很近。
\begin{itemize}
	\item 建議走法：於辛亥路進入，於萬芳交流道離開；可以在木柵動物園前的新光路稍作休息後原路返回。
	\item 建議速度：高速，單純簡單方便的兜風。
	\item 推薦指數：★☆☆☆☆
\end{itemize}
\subsection*{路線四：敦化南北路}
臺北市著名官道。夜晚車少的時候開在快車道上，沿線的貓眼石反射出像跑道燈一般的引導線；駕駛感受實屬享受。
\begin{itemize}
	\item 建議走法：基隆路進入快車道，一路開到松山機場，機場內繞一圈後原路折返。
	\item 建議速度：低速，優雅從容的行駛。
	\item 推薦指數：★★★★☆
\end{itemize}
\subsection*{路線五：仁愛路}
臺北市著名官道。沿途行經臺北市內的豪宅區，路幅最寬處 100m 的超寬敞林蔭大道。過了午夜十二點後行駛在公車專用道上，與睡著的城市一起重新獲得活力。
\begin{itemize}
	\item 建議走法：仁愛圓換進入，向西行駛於路中央。
	\item 建議速度：低速，優雅從容的行駛。
	\item 推薦指數：★★★★☆
\end{itemize}
\section*{四、結論}
也許有些每天通勤臺北市的人會覺得這是通篇屁話；這是因為你沒有在對的時間上路，你沒有在對的時間體會這幾條路、體會兜風這件事、體會這座城市；挑個沒事的晚上、坐上駕駛座，試著感受兜風的樂趣吧！
%\begin{itemize}
%	\item 兜風可以讓心靈平靜；雖然有些路段建議高速行駛仍然不建過度超速。也許有些人要喜歡像《灣岸競速》一樣在高速中獲得心靈的慰藉，但我更偏好在法定速限下的漫遊，畢竟這是兜風不是飆車。
%\end{itemize}
\end{CJK}
\end{document}









































